\documentclass[11pt,a4paper]{article}

\usepackage[T1]{fontenc}
\usepackage[utf8]{inputenc}
\usepackage{lmodern}
\usepackage{microtype}
\usepackage[a4paper,margin=23mm,headheight=15pt]{geometry}
\usepackage{xcolor}
\usepackage{booktabs,longtable,tabularx,array}
\usepackage{enumitem}
\usepackage{fancyhdr}
\usepackage{hyperref}
\usepackage{url}
\usepackage{listings}
\usepackage[most]{tcolorbox}
\usepackage{titlesec}
\usepackage{lastpage}
\usepackage{pifont}
\usepackage{amssymb}
\usepackage{seqsplit}

\definecolor{AionNavy}{HTML}{13233A}
\definecolor{AionBlue}{HTML}{245A8D}
\definecolor{AionLightBlue}{HTML}{EAF5FC}
\definecolor{AionTeal}{HTML}{168C8C}
\definecolor{AionGreen}{HTML}{237A57}
\definecolor{AionAmber}{HTML}{B97618}
\definecolor{AionRed}{HTML}{A63D40}
\definecolor{AionInk}{HTML}{1F2933}
\definecolor{AionGray}{HTML}{667788}
\definecolor{AionPale}{HTML}{F3F7FA}
\definecolor{AionRule}{HTML}{D8E2EA}

\hypersetup{
  colorlinks=true,
  linkcolor=AionBlue,
  urlcolor=AionTeal,
  citecolor=AionBlue,
  pdftitle={AION Third-Party Integration, Training and Pilot Execution Standard},
  pdfauthor={AION / Tessaris / COMDEX},
  pdfsubject={Repeatable standard for integrating external business tools into AION},
  pdfkeywords={AION, Tessaris, Pilot, Vault, connectors, OAuth, MCP, tools, skills, governance}
}

\pagestyle{fancy}
\fancyhf{}
\lhead{\textcolor{AionGray}{AION Third-Party Integration Standard}}
\rhead{\textcolor{AionGray}{Version 1.0 --- 25 August 2026}}
\cfoot{\textcolor{AionGray}{\thepage\ / \pageref{LastPage}}}
\renewcommand{\headrulewidth}{0.4pt}
\renewcommand{\headrule}{\hbox to\headwidth{\color{AionRule}\leaders\hrule height \headrulewidth\hfill}}

\titleformat{\section}{\Large\bfseries\color{AionNavy}}{\thesection}{0.7em}{}
\titleformat{\subsection}{\large\bfseries\color{AionBlue}}{\thesubsection}{0.7em}{}
\titleformat{\subsubsection}{\normalsize\bfseries\color{AionInk}}{\thesubsubsection}{0.7em}{}
\titlespacing*{\section}{0pt}{2.3ex plus .4ex}{1.0ex}
\titlespacing*{\subsection}{0pt}{1.8ex plus .3ex}{0.7ex}

\setlist[itemize]{leftmargin=1.5em,itemsep=0.25em,topsep=0.4em}
\setlist[enumerate]{leftmargin=1.7em,itemsep=0.35em,topsep=0.4em}
\renewcommand{\arraystretch}{1.18}
\setlength{\tabcolsep}{4pt}
\setlength{\parindent}{0pt}
\setlength{\parskip}{0.55em}

\newcommand{\code}[1]{\texttt{#1}}
\newcommand{\pass}{\textcolor{AionGreen}{\ding{51}}}
\newcommand{\fail}{\textcolor{AionRed}{\ding{55}}}
\newcommand{\must}{\textbf{MUST}}
\newcommand{\mustnot}{\textbf{MUST NOT}}
\newcommand{\should}{\textbf{SHOULD}}

\newtcolorbox{rulebox}{
  colback=AionBlue!4,colframe=AionBlue!75!black,
  title=Governing rule,fonttitle=\bfseries,
  boxrule=0.7pt,arc=2pt,left=6pt,right=6pt,top=5pt,bottom=5pt
}
\newtcolorbox{truthbox}{
  colback=AionRed!4,colframe=AionRed!70!black,
  title=Truth and claim boundary,fonttitle=\bfseries,
  boxrule=0.7pt,arc=2pt,left=6pt,right=6pt,top=5pt,bottom=5pt
}
\newtcolorbox{releasebox}{
  colback=AionGreen!5,colframe=AionGreen!75!black,
  title=Release condition,fonttitle=\bfseries,
  boxrule=0.7pt,arc=2pt,left=6pt,right=6pt,top=5pt,bottom=5pt
}
\newtcolorbox{warningbox}{
  colback=AionAmber!6,colframe=AionAmber!80!black,
  title=Fail-closed requirement,fonttitle=\bfseries,
  boxrule=0.7pt,arc=2pt,left=6pt,right=6pt,top=5pt,bottom=5pt
}

\lstdefinelanguage{AionJSON}{
  basicstyle=\ttfamily\footnotesize,
  string=[s]{\"}{\"},
  stringstyle=\color{AionTeal},
  comment=[l]{//},
  commentstyle=\color{AionGray}\itshape,
  morekeywords={true,false,null},
  keywordstyle=\color{AionBlue}\bfseries
}
\lstset{
  language=AionJSON,
  basicstyle=\ttfamily\footnotesize,
  backgroundcolor=\color{AionPale},
  frame=single,
  rulecolor=\color{AionRule},
  breaklines=true,
  breakatwhitespace=false,
  columns=fullflexible,
  keepspaces=true,
  showstringspaces=false,
  xleftmargin=0.5em,xrightmargin=0.5em,
  aboveskip=0.8em,belowskip=0.8em
}

\begin{document}

\begin{titlepage}
  \pagecolor{AionNavy}
  \color{white}
  \vspace*{16mm}
  {\fontsize{23}{28}\selectfont\bfseries AION Third-Party Integration, Training and Pilot Execution Standard\par}
  \vspace{4mm}
  {\fontsize{15}{19}\selectfont One repeatable method for connecting, teaching, governing and operating every external business tool\par}
  \vspace{14mm}
  {\color{AionTeal!45!white}\rule{\textwidth}{1.4pt}}
  \vspace{10mm}

  \begin{tabularx}{\textwidth}{>{\bfseries}p{39mm}X}
    Document date & 25 August 2026 \\
    Version & 1.0 \\
    System & Tessaris / AION Central Pilot \\
    Repository & \path{/Users/kevinrobinson/dev/COMDEX} \\
    Applies to & OAuth, API-key, MCP, webhook, local-service, browser and desktop-control integrations \\
    Primary outcome & Every installed provider behaves consistently, truthfully and safely from Vault connection to verified Pilot execution \\
    Status & Normative implementation and release standard \\
  \end{tabularx}

  \vfill
  \begin{tcolorbox}[colback=AionBlue!55!AionNavy,colframe=AionTeal!70!white,
    title=Purpose,fonttitle=\bfseries,coltitle=white,coltext=white,boxrule=0.6pt]
  This standard defines how any new third party becomes a genuine AION
  capability. It covers provider selection, credential storage, machine-readable
  action contracts, AION training, Pilot routing, owner authority, execution,
  read-back verification, user-facing capability explanations, failure handling,
  revocation and release evidence. The standard prevents a connected logo, a
  learned description or a generated draft from being mistaken for verified
  operational power.
  \end{tcolorbox}
  \vspace{5mm}
  {\small Internal engineering standard --- local-first, evidence-backed, approval-governed}
\end{titlepage}
\nopagecolor
\color{AionInk}

\pagenumbering{roman}
\tableofcontents
\clearpage

\section*{Document control}
\addcontentsline{toc}{section}{Document control}

\begin{tabularx}{\textwidth}{>{\bfseries}p{38mm}X}
Owner & Tessaris Connectors and AION Runtime \\
Review authority & Product owner, connector maintainer and security reviewer \\
Change rule & Changes to actions, scopes, authority, receipts or release gates require a new manifest hash and regression evidence \\
Source standard & \path{docs/aion/AION_THIRD_PARTY_TOOL_INTEGRATION_STANDARD_2026-08-24.md} \\
Machine validator & \path{backend/modules/connectors/tool_manifest.py} \\
Provider catalogue & \path{backend/modules/connectors/priority_provider_catalog.py} \\
Release gate & \path{backend/modules/connectors/provider_release_gate.py} \\
Human checklist & \path{docs/aion/templates/AION_THIRD_PARTY_TOOL_IMPLEMENTATION_CHECKLIST.md} \\
\end{tabularx}

\begin{truthbox}
An integration is not operational merely because AION knows how the provider
works, a credential has been entered, an adapter exists, or a dry run passes.
Live capability is action-specific and exists only when the current workspace
has a healthy grant, sufficient scope, applicable owner authority, a released
adapter and the required verification receipts. The Vault and Pilot must report
that state exactly.
\end{truthbox}

\clearpage
\pagenumbering{arabic}

\section{Executive standard}

Every third party is represented as a collection of narrow, named powers. Gmail
is not one power called ``use Gmail.'' It is a set such as search email, read
email, create draft, send email, reply, label and archive. Each power has its own
inputs, risk, permission scope, owner authority, adapter method, receipt and
release evidence.

The definitive capability statement is computed from five sources:

\begin{enumerate}
  \item the validated provider manifest;
  \item the provider's achieved release stage;
  \item the current workspace credential and granted scopes;
  \item the current user's and workspace's authority policy; and
  \item current connector health and incident state.
\end{enumerate}

The same computed statement drives the Vault connection screen, Pilot's action
resolver, the answer to ``what can you do with this tool?'', approval screens
and the execution gate. No surface may maintain a separate hard-coded list.

\begin{rulebox}
\textbf{One manifest, one live-state resolver, many views.} The Vault explains
the powers; Pilot selects and executes them; the dashboard displays their state;
and receipts prove their outcomes. All four read the same action IDs and current
eligibility. This eliminates drift between what the user is promised and what
the runtime can actually do.
\end{rulebox}

\section{Scope and terminology}

\begin{longtable}{>{\bfseries}p{34mm}p{118mm}}
\toprule
Term & Meaning \\
\midrule
Provider & The external system, for example Gmail, Microsoft Outlook, Xero, HubSpot, Twilio or Stripe. \\
Connector & The authenticated channel between Tessaris and the provider. \\
Adapter & Code that converts a validated AION action into a bounded provider operation and reconciles the result. \\
Action & A stable atomic power such as \code{gmail.send\_email}. \\
Skill & AION's learned method for deciding when and how an action should be used. \\
Business knowledge & Owner-approved company facts, policies and operating preferences retrieved from the Business Brain. \\
Vault reference & An opaque identifier for locally stored credentials; never the secret itself. \\
Authority & The owner's permission mode and limits for a named action. \\
Receipt & Durable local evidence of the requested payload, policy decision, provider result and reconciliation. \\
Read-back & A separate provider query that confirms the resulting external state. \\
Capability state & A runtime result describing whether an action is known, connected, executable, approval-gated or blocked. \\
Release stage & The highest evidence-backed maturity level achieved by the provider integration. \\
\bottomrule
\end{longtable}

\section{Reference architecture}

\begin{tcolorbox}[colback=AionPale,colframe=AionRule,boxrule=0.6pt]
\begin{verbatim}
User request
    |
    v
AION intent + Business Brain + prior experience
    |
    v
Capability resolver ---- provider manifest
    |                     release receipts
    |                     Vault health/scopes
    |                     owner authority
    v
Pilot plan and exact preview
    |
    v
Authority decision ---- blocked / ask / bounded auto / full access
    |
    v
Adapter prepare -> execute -> provider read-back
    |
    v
Receipt + business outcome + experience capsule
\end{verbatim}
\end{tcolorbox}

The language model never receives raw credentials. It proposes a stable action
and structured payload. A deterministic control plane decides whether that
action exists and is currently eligible. The adapter receives a short-lived
Vault reference only after policy evaluation.

\subsection{Separation of responsibilities}

\begin{tabularx}{\textwidth}{>{\bfseries}p{32mm}X X}
\toprule
Layer & Responsible for & Must not do \\
\midrule
AION intelligence & Understand intent, retrieve approved knowledge, choose a skill, fill or request missing inputs & See secrets, invent permissions or claim execution \\
Capability resolver & Join manifest, release, connection, scope, authority and health state & Make provider calls \\
Pilot & Present the plan, obtain approval when required, dispatch eligible work and explain results & Bypass policy or hard-code provider behaviour \\
Vault & Store credentials, scopes, connection metadata and revocation controls & Expose raw secrets to prompts or receipts \\
Adapter & Validate, preview, execute and reconcile a named action & Accept arbitrary provider commands \\
Receipt ledger & Preserve exact evidence and uncertainty & Convert an attempted request into unverified success \\
\bottomrule
\end{tabularx}

\section{The standard provider installation lifecycle}

Every integration \must follow the stages below in order. A stage may contain
several actions, but no later claim can be made from evidence belonging only to
an earlier stage.

\subsection{Stage 1: select the business outcome}

Define the useful job before selecting technology. Examples include reading an
inbox, preparing and sending an approved reply, creating a CRM contact, issuing
an invoice, placing a consented call or reconciling a payment. Record the owning
department, expected user value and prohibited uses.

\subsection{Stage 2: decompose the provider into atomic actions}

Each operation receives a stable provider-prefixed action ID. Read, draft,
external write, delete and financial operations remain distinct. A universal
``do anything'' action is prohibited.

\subsection{Stage 3: choose and register the connection}

Supported connection types are OAuth~2, scoped API key, signed webhook, MCP,
local service, browser session and desktop control. Official APIs or official
MCP servers are preferred when they expose the required result identifiers and
read-back operations. Browser or desktop control is a bounded fallback, not a
universal permission.

The connection definition \must include:

\begin{itemize}
  \item a Vault reference;
  \item least-privilege scopes;
  \item a side-effect-free health check;
  \item token-expiry and refresh behaviour;
  \item a revocation method and emergency kill switch; and
  \item confirmation that secrets are invisible to the model, logs and receipts.
\end{itemize}

\subsection{Stage 4: create and validate the tool manifest}

The manifest is the provider's machine-readable contract. It declares identity,
connection, required skills, required company knowledge, actions, risks,
authority, adapter methods, receipts, reconciliation and failure behaviour. It
is validated and SHA-256 fingerprinted before registration.

\subsection{Stage 5: train and examine AION}

AION learns the provider model and each action's correct use. Training does not
grant credentials or authority. Section~\ref{sec:training} defines the full
curriculum and evidence gates.

\subsection{Stage 6: implement the adapter}

The adapter implements \code{health\_check}, \code{prepare}, \code{preview},
\code{execute}, \code{reconcile} and \code{revoke} where applicable. Writes use
idempotency and fail closed. The adapter accepts only registered actions and
validated payloads.

\subsection{Stage 7: bind authority}

Every action is bound to an owner policy: \code{ask\_each\_time},
\code{auto\_within\_limits}, \code{full\_access} or \code{blocked}. The policy
may add recipients, accounts, records, hours, amounts, currencies or frequency
limits. A refusal with \code{ask\_again=false} is retained.

\subsection{Stage 8: verify progressively}

Contract tests and dry runs come first. A sandbox action and provider read-back
come next. Finally, one bounded live action receives exact approval, executes,
is read back from the provider, and produces durable receipts. Failure and
revocation must also be proved.

\subsection{Stage 9: release and monitor}

The action is displayed as live only at its achieved stage. Runtime health,
scope loss, provider incidents or revoked authority can immediately demote or
block it without editing the manifest.

\section{The machine-readable provider contract}

The canonical schema is \code{aion.third\_party\_tool.v1}. The repository
template is:

\begin{center}
\path{backend/modules/connectors/templates/provider_tool_manifest.template.json}
\end{center}

The following abbreviated example shows the required shape. It is illustrative;
the validated repository manifest remains authoritative.

\begin{lstlisting}
{
  "schema_version": "aion.third_party_tool.v1",
  "tool": {
    "tool_id": "gmail",
    "label": "Gmail",
    "category": "email"
  },
  "connection": {
    "type": "oauth2",
    "vault_reference": "vault.gmail.credentials",
    "least_privilege_scopes": [
      "gmail.readonly", "gmail.compose", "gmail.modify"
    ],
    "health_check": "gmail.health_check",
    "revocation_method": "revoke grant and remove Vault entry"
  },
  "skills": [{
    "skill_id": "gmail.business_email",
    "required_knowledge": [
      "email_etiquette", "business_identity",
      "recipient_policy", "approved_company_knowledge"
    ]
  }],
  "actions": [{
    "action_id": "gmail.send_email",
    "kind": "write",
    "risk": {"external_write": true, "tier": "medium"},
    "authority": {
      "default_mode": "ask_each_time",
      "allowed_modes": [
        "ask_each_time", "auto_within_limits",
        "full_access", "blocked"
      ]
    },
    "execution": {
      "adapter_method": "send_email",
      "dry_run_supported": true,
      "idempotency": "required"
    },
    "receipt": {"required": true},
    "failure": {"fail_closed": true}
  }]
}
\end{lstlisting}

\subsection{Mandatory action fields}

\begin{longtable}{>{\bfseries}p{34mm}p{118mm}}
\toprule
Field & Required content \\
\midrule
Action identity & Stable action ID, user label and plain-English purpose. \\
Skills & Skill IDs and approved business knowledge required before use. \\
Input contract & Required fields, types, formats, limits and missing-input behaviour. \\
Output contract & Provider request IDs, result IDs, status and completion evidence. \\
Risk & Read/write, financial status, reversibility, data classes and tier. \\
Authority & Policy key, default mode, allowed modes and constraints. \\
Execution & Adapter method, dry-run support, timeout and idempotency. \\
Receipt & Required evidence fields and reconciliation method. \\
Failure & Fail-closed rule, retryable classes and truthful user message. \\
Explanation & User-facing summary of what the action does and what it does not do. \\
\bottomrule
\end{longtable}

\section{Vault connection and capability explanation}
\label{sec:vault}

The Vault is both the secure connection point and the user's capability
contract. Before the user connects a provider, the screen explains what AION
could do if the required scopes and adapters are available. After connection,
it explains what AION can actually do in the current workspace.

\subsection{Required provider card}

Every provider card \must show:

\begin{itemize}
  \item provider name, category and connection method;
  \item connection health and last successful check;
  \item requested permissions in plain English;
  \item \textbf{AION can now}: currently eligible read, draft and write actions;
  \item \textbf{AION will ask before}: approval-gated actions;
  \item \textbf{AION cannot yet}: missing scopes, adapters, authority or release evidence;
  \item current authority mode and important limits;
  \item data accessed and local evidence retained;
  \item buttons to connect, test, change authority and revoke; and
  \item a link to the detailed provider capability record.
\end{itemize}

The screen \mustnot display ``connected'' as a synonym for ``can execute.'' It
must separately display authentication, scope, adapter, release and authority
state.

\subsection{Example: Gmail before connection}

\begin{tcolorbox}[colback=AionLightBlue,colframe=AionBlue!45]
\textbf{Gmail --- business email}\par
Connect Gmail to let AION search and read permitted mail, prepare drafts, and
send or reply when the relevant action is verified and authorised. AION never
sees your password or OAuth token. Connecting Gmail does not automatically
permit sending, deleting or acting on every message.
\end{tcolorbox}

\subsection{Example: Gmail after connection}

\begin{tcolorbox}[colback=AionLightBlue,colframe=AionBlue!45]
\textbf{Connected --- healthy grant}\par
\textbf{AION can now:} search permitted email, read messages and prepare drafts.\par
\textbf{AION will ask before:} sending a new email or reply under the current
\code{ask\_each\_time} policy.\par
\textbf{AION cannot:} send without the required scope and released adapter,
read excluded accounts, reveal credentials, or claim delivery without provider
evidence.\par
\textbf{Permissions:} read, compose and modify mail. \textbf{Last checked:}
derived from the live connector health receipt.
\end{tcolorbox}

\subsection{Capability explanation must be generated, not written twice}

The Vault UI does not own these sentences. A shared capability-introspection
service renders them from the current computed state. Pilot uses the same
service when answering questions about tools. Provider labels may have curated
plain-English wording in the manifest, but the verbs and eligibility state come
from registered actions and live evidence.

\section{Training AION on a provider}
\label{sec:training}

AION must learn more than the provider's marketing description or API syntax.
It must learn how business objects behave, which operation achieves which
outcome, what can go wrong and how to prove the result.

\subsection{Seven mandatory lessons}

\begin{enumerate}
  \item provider objects, identifiers and state transitions;
  \item least-privilege authentication and local Vault handling;
  \item atomic read, draft and write action contracts;
  \item business policy, authority limits and approval escalation;
  \item provider errors, retries, idempotency and rate limits;
  \item receipt creation and independent provider reconciliation; and
  \item credential revocation, kill switch and incident response.
\end{enumerate}

Training sources should include official provider documentation, the actual
adapter contract, workspace policy, approved Business Brain facts, representative
fixtures and observed failure receipts. Community tutorials may support
discovery but do not override official or locally verified evidence.

\subsection{Training outputs}

The training run produces:

\begin{itemize}
  \item a provider object and state model;
  \item an action-selection map from user intent to stable action IDs;
  \item required and optional input definitions;
  \item permission and authority explanations;
  \item error, retry and reconciliation playbooks;
  \item adversarial examples and prohibited assumptions;
  \item verified experience capsules from practical tests; and
  \item a source and version ledger so learning can be refreshed when the provider changes.
\end{itemize}

\subsection{Required assessments}

\begin{tabularx}{\textwidth}{>{\bfseries}p{44mm}X}
\toprule
Assessment & Evidence required \\
\midrule
Closed-book provider model & AION explains objects, states, actions, scopes and failure cases without source access. \\
Intent-to-action transfer & Unfamiliar user requests map to the correct action or to an honest capability gap. \\
Safe dry run & AION constructs a valid payload and preview with zero external side effects. \\
Adversarial permission test & AION refuses scope escalation, hidden recipients, unsafe amounts and unsupported actions. \\
Sandbox execution & A controlled write completes and is independently read back. \\
Failure and repair & AION diagnoses expired credentials, rate limits, duplicates and ambiguous outcomes without inventing success. \\
Retention & Closed-book checks after 1, 7, 30 and 90 days, with source-free execution where precommitted. \\
\bottomrule
\end{tabularx}

\begin{warningbox}
Reading documentation is not provider competence. Passing a knowledge test is
not operational readiness. AION is classified as practically competent only
after unfamiliar actions, failures, repairs and provider read-backs produce
independent evidence. Training never supplies credentials or owner authority.
\end{warningbox}

\section{Making Pilot genuinely able to do the work}

Pilot must route a natural request to installed powers rather than match it to
a hard-coded response. The same mechanism must handle ``send that Gmail
message,'' ``raise this Xero invoice,'' ``move the HubSpot deal'' and future
providers without adding one-off conversational branches.

\subsection{Runtime decision sequence}

For every user request Pilot performs the following sequence:

\begin{enumerate}
  \item identify the intended business outcome and affected objects;
  \item retrieve approved Business Brain facts, policies and relevant experience;
  \item query the capability resolver for matching action IDs;
  \item evaluate release stage, connection, scope, health and authority;
  \item select an eligible action, a safe draft alternative or a precise blocker;
  \item ask only for information genuinely missing from the action contract;
  \item produce the exact preview and payload hash;
  \item execute only after the applicable authority decision;
  \item reconcile provider state independently;
  \item record the receipt, business outcome and any learned repair; and
  \item answer the user in plain English with what happened and what remains uncertain.
\end{enumerate}

\subsection{Capability resolver contract}

The resolver should expose a function equivalent to:

\begin{lstlisting}
describe_provider_capabilities(
  workspace_id,
  user_id,
  provider_id,
  optional_intent
) -> ProviderCapabilityView
\end{lstlisting}

Its output must be deterministic from current state and suitable for both UI
and runtime use:

\begin{lstlisting}
{
  "provider_id": "gmail",
  "connection": "healthy",
  "release_stage": "live_verified",
  "actions": [
    {
      "action_id": "gmail.search_email",
      "status": "available",
      "authority": "auto_within_limits"
    },
    {
      "action_id": "gmail.create_draft",
      "status": "available",
      "authority": "auto_within_limits"
    },
    {
      "action_id": "gmail.send_email",
      "status": "approval_required",
      "authority": "ask_each_time"
    }
  ],
  "blocked": [],
  "generated_from": {
    "manifest_sha256": "...",
    "credential_health_receipt": "...",
    "authority_revision": "..."
  }
}
\end{lstlisting}

\subsection{Answering ``What can you do with Gmail?''}

Pilot first queries the resolver, then gives a state-qualified answer. For
example:

\begin{tcolorbox}[colback=AionPale,colframe=AionRule]
\textbf{Connected and approval-gated}\par
``With the Gmail account connected to this business, I can search and read the
permitted inbox and prepare email drafts. I can also send new emails and replies,
but your current policy requires me to show you the exact recipient, subject and
message for approval first. I cannot access another account, see your OAuth
secret, or confirm delivery unless Gmail provides verifiable status.''
\end{tcolorbox}

If disconnected, Pilot explains the available integration and offers to open
the Vault. If a scope is missing, it names the affected actions. If the adapter
is not released, it says that the action is not yet operational. These are
computed answers, not provider-specific dialogue scripts.

\subsection{Missing information behaviour}

Pilot should request the minimum missing input conversationally. A draft email
may proceed with clearly marked placeholders; an external send cannot. The
action contract determines whether Pilot asks, uses an approved default,
prepares a placeholder draft or stops. It must not manufacture recipients,
amounts, legal commitments, account identifiers or company facts.

\section{Authority and removal of the stabilisers}

The owner can increase autonomy without changing the provider adapter. Authority
is action-specific, visible, reversible and bounded.

\begin{longtable}{>{\bfseries}p{38mm}p{113mm}}
\toprule
Mode & Behaviour \\
\midrule
\code{blocked} & Never execute. Optionally retain \code{ask\_again=false}. \\
\code{ask\_each\_time} & Show the exact payload and request approval for this instance. \\
\code{auto\_within\_limits} & Execute only within declared recipients, accounts, amounts, times, frequency and other constraints. \\
\code{full\_access} & Standing authority for the named action, still subject to adapter validation, provider scope, kill switch, receipts and statutory safety constraints. \\
\bottomrule
\end{longtable}

Read authority never implies write authority. Draft authority never implies
send authority. Invoice creation never implies issue, payment or bank transfer.
High-risk actions cannot default to unrestricted full access. Financial actions
must define amount and currency controls, and may impose a mandatory approval
floor regardless of the general mode.

When authority is missing, Pilot may ask once whether the owner wants to approve
the instance or change the standing rule. It must explain the exact capability,
scope and consequences. If the owner declines and selects ``do not ask again,''
Pilot records the blocked policy and later provides a quiet route to manual
settings instead of repeatedly prompting.

\section{Adapter execution contract}

\subsection{Required phases}

\begin{enumerate}
  \item \code{health\_check()} proves the connection without changing provider state;
  \item \code{prepare()} validates identities, inputs, policy and provider constraints;
  \item \code{preview()} renders exactly what is proposed;
  \item \code{execute()} performs only the named registered action;
  \item \code{reconcile()} reads provider state back and binds it to the payload hash; and
  \item \code{revoke()} or the central kill switch prevents further execution.
\end{enumerate}

\subsection{Write safety}

External writes must use an idempotency key or verified provider-managed
idempotency. Timeouts cannot be treated as failure or success until reconciled.
Retries cannot duplicate messages, invoices, calls or payments. Arbitrary raw
API calls remain blocked unless deliberately registered as a separate bounded
action.

\section{Receipts, read-back and learning from outcomes}

Every consequential action produces a local receipt containing:

\begin{itemize}
  \item workspace, provider, action and adapter version;
  \item manifest hash and authority decision;
  \item exact payload hash, without secrets;
  \item provider request and result identifiers;
  \item attempt and completion timestamps;
  \item reconciliation method and result;
  \item final state: \code{verified\_complete}, \code{failed} or \code{uncertain}; and
  \item error classification and safe repair when applicable.
\end{itemize}

Successful and failed practical work becomes a governed AION experience
capsule. Future retrieval may reuse the verified method and repair, but never a
credential, unverified claim or workspace-private fact outside its authority.

\section{Release stages and exact evidence gate}

\begin{longtable}{>{\bfseries}p{36mm}p{115mm}}
\toprule
Stage & Meaning \\
\midrule
\code{catalogued} & Provider and intended business value are registered. \\
\code{contract\_ready} & Valid manifest and contract SHA-256 exist. \\
\code{knowledge\_ready} & AION passed the provider knowledge assessment. \\
\code{adapter\_ready} & Adapter contract and dry-run receipts pass. \\
\code{credential\_ready} & Current workspace credential health is verified. \\
\code{sandbox\_verified} & Controlled provider write and read-back pass. \\
\code{live\_verified} & Approved live action, independent read-back, failure test and revocation test all pass. \\
\bottomrule
\end{longtable}

The machine gate requires the following evidence before live verification:

\begin{enumerate}
  \item \code{contract\_sha256};
  \item \code{knowledge\_assessment\_receipt};
  \item \code{adapter\_contract\_receipt};
  \item \code{dry\_run\_receipt};
  \item \code{credential\_health\_receipt};
  \item \code{sandbox\_write\_receipt};
  \item \code{approved\_live\_action\_receipt};
  \item \code{provider\_readback\_receipt};
  \item \code{failure\_test\_receipt}; and
  \item \code{revocation\_test\_receipt}.
\end{enumerate}

\begin{releasebox}
Only \code{live\_verified} permits the statement that AION can perform the
tested action through the provider. The statement remains qualified by current
workspace connection, scopes, health and authority. Release is per action where
provider actions mature at different rates.
\end{releasebox}

\section{End-to-end Gmail example}

This example is the reference pattern for future integrations.

\subsection{Business scope}

The initial Gmail release supports searching permitted mail, reading messages,
creating drafts, sending new messages and replying. Delete, account-wide
administration and unrestricted mailbox automation are excluded until separately
contracted.

\subsection{Connection and Vault}

The workspace completes OAuth, granting only the scopes required by enabled
actions. The Vault stores tokens locally under \code{vault.gmail.credentials}.
The model sees only \code{gmail} capability state and never the token. The Vault
tests connection health and offers immediate revocation.

\subsection{Training}

AION learns Gmail message/thread identifiers, drafts, replies, headers, recipient
handling, scopes, common API failures, retry/idempotency constraints and read-back.
It also retrieves business identity, approved company facts, email etiquette and
recipient policy from the Business Brain.

\subsection{Pilot request}

For ``reply to Jane confirming Friday's installation,'' Pilot:

\begin{enumerate}
  \item resolves the relevant thread and \code{gmail.reply\_email};
  \item confirms Jane's identity and the approved installation fact;
  \item drafts the reply without requiring approval;
  \item checks current send authority;
  \item if required, shows the exact account, recipient, subject and body;
  \item executes once approved or automatically within a standing bounded rule;
  \item reads back the sent-message identifier and thread state; and
  \item reports verified completion or an honest uncertain/failed state.
\end{enumerate}

\subsection{Capability answer}

When asked what Gmail can do, Pilot uses the resolver. It does not simply recite
the manifest. Its answer changes if the grant expires, sending scope is absent,
the action is blocked, or a provider incident demotes the connector.

\section{Testing requirements}

Every provider release must demonstrate all of the following:

\begin{itemize}
  \item manifest validation and stable hash;
  \item secrets absent from manifest, prompt, log and receipt;
  \item missing, expired and revoked credential handling;
  \item insufficient-scope handling;
  \item deterministic capability explanations in Vault and Pilot;
  \item proof that UI wording changes when live state changes;
  \item intent selection without provider-specific hard-coded conversation logic;
  \item safe drafts without unnecessary approval;
  \item blocked and approval-required writes;
  \item exact-payload preview and hash binding;
  \item zero-side-effect dry run;
  \item idempotent duplicate and timeout behaviour;
  \item rate-limit and malformed-response handling;
  \item sandbox execution and provider read-back;
  \item one bounded approved live action and independent read-back;
  \item failure, kill-switch and revocation tests; and
  \item no claim of completion when outcome is uncertain.
\end{itemize}

\section{Standard implementation checklist}

\subsection*{Provider and contract}
\begin{itemize}
  \item[$\square$] Business outcome, owner and prohibited uses recorded.
  \item[$\square$] Atomic actions and user explanations defined.
  \item[$\square$] Manifest validates and SHA-256 is preserved.
  \item[$\square$] Official provider documentation version is recorded.
\end{itemize}

\subsection*{Connection and Vault}
\begin{itemize}
  \item[$\square$] Least-privilege connection and scopes implemented.
  \item[$\square$] Credentials remain local and invisible to the model.
  \item[$\square$] Health check, token refresh, revocation and kill switch work.
  \item[$\square$] Vault card displays actual powers, approvals, blockers and limits.
\end{itemize}

\subsection*{Training and runtime}
\begin{itemize}
  \item[$\square$] Seven-lesson provider syllabus completed.
  \item[$\square$] Closed-book, transfer and adversarial assessments pass.
  \item[$\square$] Pilot resolves natural requests to stable action IDs.
  \item[$\square$] Capability answers come from the live resolver.
  \item[$\square$] Missing inputs are requested naturally and minimally.
\end{itemize}

\subsection*{Execution and release}
\begin{itemize}
  \item[$\square$] Adapter prepare, preview, execute, reconcile and revoke paths work.
  \item[$\square$] Authority modes and limits are visible and enforced.
  \item[$\square$] Writes are idempotent and receipt-backed.
  \item[$\square$] Sandbox and bounded live actions are independently read back.
  \item[$\square$] Failure and revocation evidence pass the machine release gate.
  \item[$\square$] The published claim matches the achieved action stage exactly.
\end{itemize}

\section{Operational maintenance}

Provider contracts require continuing maintenance. Scope names, API versions,
rate limits, schemas and consent rules can change. A scheduled health process
should detect documentation or provider-version changes, rerun contract tests,
and demote affected actions until refreshed evidence passes. Incidents must be
visible in the Vault and Pilot capability answer.

New community skills enter quarantine. Their source, licence, dependencies,
network destinations, secret access and command surface are reviewed before
installation. They never inherit business credentials or standing authority.

\section{Current implementation map}

\begin{longtable}{>{\bfseries}p{48mm}p{103mm}}
\toprule
Purpose & Repository path \\
\midrule
Manifest validator & \path{backend/modules/connectors/tool_manifest.py} \\
Manifest template & \path{backend/modules/connectors/templates/provider_tool_manifest.template.json} \\
Priority provider manifests and curricula & \path{backend/modules/connectors/priority_provider_catalog.py} \\
Release evidence gate & \path{backend/modules/connectors/provider_release_gate.py} \\
Curated AION business skill catalogue & \path{backend/modules/connectors/aion_business_skill_catalog.py} \\
Manifest tests & \path{backend/tests/connectors/test_tool_manifest.py} \\
Provider and gate tests & \path{backend/tests/connectors/test_priority_provider_catalog.py} \\
Human provider checklist & \path{docs/aion/templates/AION_THIRD_PARTY_TOOL_IMPLEMENTATION_CHECKLIST.md} \\
Concise integration standard & \path{docs/aion/AION_THIRD_PARTY_TOOL_INTEGRATION_STANDARD_2026-08-24.md} \\
Provider expansion programme & \path{docs/aion/AION_PROVIDER_AND_SKILL_EXPANSION_PROGRAMME_2026-08-24.md} \\
\bottomrule
\end{longtable}

\section{Current provider claim boundary}

The repository presently contains validated contracts and training syllabi for
Gmail, Microsoft Outlook, Xero, HubSpot, Twilio Voice and Stripe. Repository
readiness records describe foundations and blockers, not universal live power.
At the time of this document, no provider has evidence for every
\code{live\_verified} release field in the central gate.

\begin{truthbox}
The standard is complete as a repeatable engineering contract. It does not
claim that every provider, every account or every action is live. AION may say
``I can'' only when the requested action is registered, released, connected,
healthy, sufficiently scoped and authorised in the current workspace. Otherwise
it must say what it can prepare safely, what is missing and how the owner can
enable or approve the capability.
\end{truthbox}

\section{Definition of a successful installation}

A new third party is successfully installed when a user can connect it in the
Vault, understand its real powers before granting access, ask Pilot an unfamiliar
natural-language question about those powers, receive a truthful state-qualified
answer, delegate an eligible task, review or automatically authorise it under a
visible policy, obtain a verified provider result and receipt, revoke access, and
observe Pilot immediately stop claiming or exercising that power.

That end-to-end behaviour---not the existence of a logo, credential, prompt,
skill file or adapter method---is the Tessaris definition of a working AION
integration.

\end{document}
